\subsection{Irreducibility}
%%%%%%%%%%%%%%%%%%%%%%%%%

A crucial fact towards proving Theorem \ref{thm:fair-trees} is the irreducibility 
of $\kpw$ for elements of $\Cb$. 

\begin{proposition}\label{prop:irred-C}
If $(P,\om)$ is a connected element of $\Cb$ then $\kpw$ is irreducible in \qsym.
\end{proposition}

To prove this, we first recall the following general property.  A polynomial with integer coefficients is said to be \emph{primitive} if 
whenever an integer $k$ divides all its coefficients we have $k=\pm1$.
In the same vein, since we consider $\qsym \subseteq \mathbb{Z}[[x_1, x_2, \ldots]]$, we say $f \in \qsym$ is \emph{primitive} if whenever an integer $k$ divides $f$ in \qsym\ we have $k = \pm1$.
We are interested in whether $\kpw$ is primitive.  

To answer this question, let us define the \emph{leading term} of a formal power series $f$ expanded in terms of monomials as the term $c_\alpha x^\alpha = c_\alpha x_1^{\alpha_1} x_2^{\alpha_2} \cdots$ with lexicographically largest $\alpha$ such that $c_\alpha \neq 0$; naturally, we call this $c_\alpha$ the \emph{leading coefficient}.  For a labeled poset $(P,\om)$, we define $x^{\mathrm{jump}(P,\om)} = x_1^{j_1} x_2^{j_2} \cdots$ where $j_i$ is the number of elements of $(P,\om)$ of jump $i-1$.

\begin{proposition}[{\cite[proof of Proposition 4.2]{McWa14}}]\label{prop:prim}
For any labeled poset $(P,\om)$, the leading term of $\kpw$ is $x^{\mathrm{jump}(P,\om)}$.  In particular, the leading coefficient of $K_{(P, \om)}$ is $1$ and $K_{(P, \om)}$ is primitive.
\end{proposition}

The next lemma is relevant to the recursive construction of $\Cb$.

\begin{lemma}\label{lem:irred}
If $f$ is a primitive element of \qsym\  then the polynomials
$F_1\wP f$,\ $F_1\sP f$,\ $f\wP F_1$, and $f\sP F_1$
are irreducible in \qsym.
\end{lemma}
\begin{proof}
Let us first deal with $F_1\sP f$.
Assume that $F_1\sP f$ is reducible. 
Note that if $f$ expands as $\sum c_\alpha F_\alpha$ then $F_1\sP f = \sum c_\alpha F_{\alpha^+}$, where $\alpha^+$ is obtained from $\alpha$ by appending an entry of 1 at the start.  Thus since $f$ is primitive, $F_1\sP f$ is also primitive, so there exist non-constants $g,g'\in \qsym$, such that 
\begin{equation}\label{eq:lem-irred1}
F_1\sP f = g g'.
\end{equation}
All $\alpha$ in the $F$-support of $F_1\sP f$ satisfy $\alpha_1=1$ but by Lemma~\ref{lem:deg2}, the $F$-support of the right-hand side of~\eqref{eq:lem-irred1} contains at least one $\alpha$ with $\alpha_1 \geq 2$, a contradiction.  The case of  $f\sP F_1$ is treated similarly, using the last part of $\alpha$ instead of the first part.

Let us now consider $F_1\wP f$.
We again assume $F_1\wP f$ is reducible and now $\alpha^+$ is obtained from $\alpha$ by increasing the first part by 1.  The primitivity 
implies the existence of two non-constants $g,g'\in \qsym$, 
such that $F_1\wP f = g g'$. We apply the bar operator to this equality
and get 
\[
\inv{F_1\wP f} = \inv{g g'} = \inv{g} \inv{g'}.
\]
by Lemma~\ref{lem:bar}\ref{it47b}. Since the elements $\alpha$ of the $F$-support of $F_1\wP f$ all satisfy $\alpha_1\ge 2$, we know all $\alpha'$ in the $F$-support of 
$\inv{F_1\wP f}$ satsify $\alpha'_1 = 1$.
We are thus led to the same contradiction as in the first case. 
The case of  $f\wP F_1$ is treated similarly.
\end{proof}

Since we will only be considering labeled posets in $\Cb$ for the remainder of this section, we will abbreviate $(P,\om)$ as $P$ for easier reading.

\begin{proof}[Proof of Proposition \ref{prop:irred-C}]
If $P$ has just a single element, then $K_P = F_1$, which is irreducible.  Otherwise, since $P$ is connected and constructed recursively, it must be of one of the forms $[1]\wP P'$,\ $[1]\sP P'$,\ 
$P'\wP [1]$, or $P'\sP [1]$, where $P' \in \Cb$.  The result now follows from Propositions~\ref{prop:k_products} and~\ref{prop:prim}, and Lemma~\ref{lem:irred}.
\end{proof}

\begin{remark}
A special case of Proposition~\ref{prop:irred-C}, or even just Lemma~\ref{lem:irred}, is that $F_\alpha$ is irreducible in \qsym, as previously shown in~\cite{LaPy08}.
\end{remark}
